\section{Validation and Sensitivity Analysis}

The first validation layer is section-capacity feasibility. After the operation plan and timetable are fixed, the overlap zone and the non-overlap sections must both be checked against their effective carrying capacity. This audit is more informative than a narrative statement about “sufficient transport ability”, because it reveals whether any local section becomes the actual bottleneck of the selected service pattern.

The second validation layer concerns tracking and dwell feasibility. A timetable that looks visually regular is not automatically operationally feasible. Therefore, the minimum tracking interval between relevant train pairs and the dwell-time bounds at each stop should be checked explicitly. If these two audits are not shown separately, then the timetable layer remains only partially verified.

The third validation layer is scenario comparison. Demand growth or interval perturbation may change the service-cost tradeoff even when the baseline timetable passes all audits. For this reason, the selected service pattern should be compared under several realistic scenario settings, and the paper should state whether the final recommendation remains feasible, becomes more expensive, or requires a different short-turn intensity.

Taken together, the validation section closes the full timetable route. The selected operation plan is not judged only by its cost-service score, and the timetable is not judged only by whether it can be drawn. Instead, the final recommendation is supported jointly by capacity, tracking, dwell, and scenario evidence.
